% Page 055

\seite{55}

\abschnitt{Schuljahr 1914–18}
\pstart
Mit dem 18.\ Februar wurden 4 Kinder aus der Schule entlassen\ob
2 Knaben und zwei Mädchen, welche sich alle der Landwirtschaft

Kriegsanleihe\ob
die hiesige Schule zeichnete zur fünften Kriegsanleihe 352 M.\ob
\dittomark{die} \dittomark{hiesige} \dittomark{Schule} \dittomark{zeichnete} \dittomark{zur} 6.\ \dittomark{Kriegsanleihe} 720 M.\ob
\dittomark{die} \dittomark{hiesige} \dittomark{Schule} \dittomark{zeichnete} \dittomark{zur} 7.\ \dittomark{Kriegsanleihe} 580 \dittomark{M.}

Schuljahr 1917–18\ob
Es wurden 9 Kinder in die Schule aufgenommen\ob
4 Knaben, 5 Mädchen. Die Schülerzahl beträgt 38, davon\ob
sind 14 Knaben, 24 Mädchen.

\pend
\abschnitt{Ernte}
\pstart

Die Gesamternte war befriedigend. Aber die Getreideernte im\ob
allgemeinen sei deshalb besser weil jetzt alle mehr ar\ohb
beiteten. Heu gab's stellenweise nur wenig, so daß selbst be\ohb
mittelte Leute kaum für ihren Bedarf hatten, und vielerlei\ob
Früchte als Ersatzfrüchte anwenden mußten. Die Viehpreise\ob
waren sehr hoch. Der Preis der Lebensmittel stieg fortwährend.

\pend
\abschnitt{Schuljahr 1918–19}
\pstart

Entlassen wurden in diesem Jahre 4 Kinder.\ob
3 Mädchen und ein Knabe, sämtliche halfen\ob
bei der Landwirtschaft. Aufgenommen wurden\ob
Ostern 17.4.19 3 Kinder. die Schülerzahl ist mit 39 Schulpflichti\ohb
gen, 13 Knaben und 26 Mädchen.

\pend
\abschnitt{Ernte}
\pstart

Die Gesamternte war eine mittlere Ernte. Dem Umfange nach\ob
einen guten Ertrag. Hafer, altdeutsche Sommerfrucht, gut, sehr wenig Kartoffeln\ob
gut, auch die Viehpreise waren hoch. Der Preis der Lebensmittel stieg\ob
fortwährend.
\pend
